% Generated draft: requires mathematical and source review.
\documentclass[11pt]{article}
\usepackage[margin=25mm]{geometry}
\usepackage{fontspec}
\setmainfont{Latin Modern Roman}
\usepackage{amsmath,amssymb,mathtools,unicode-math}
\setmathfont{Latin Modern Math}
\usepackage{enumitem,needspace,xurl,hyperref}
\setlength{\parindent}{0pt}
\setlength{\parskip}{5pt}
\setlength{\emergencystretch}{3em}
\newcommand{\N}{\mathbb{N}}
\newcommand{\Z}{\mathbb{Z}}
\newcommand{\Q}{\mathbb{Q}}
\newcommand{\R}{\mathbb{R}}
\newcommand{\C}{\mathbb{C}}
\newcommand{\F}{\mathbb{F}}
\newcommand{\A}{\mathbb{A}}
\newcommand{\PP}{\mathbb{P}}
\newcommand{\E}{\mathbb{E}}
\newcommand{\eps}{\varepsilon}
\newcommand{\dd}{\,\mathrm{d}}
\begin{document}
\section*{UnsolvedMath ID: 30006270}
\textbf{Source problem code:} catalogue entry
\subsection{Definitions and mathematical statement}
The supplied record does not define a mathematical problem that can be stated precisely. It contains the following proposal, with notation transcribed as supplied:
\begin{quote}
In [4] and [5], we propose using a noncommutative space: \(T(d)w\), called a BPS category, as answer to Question 2.
\end{quote}
The parameters \(d\) and \(w\), their domains, the construction of the category, and the intended placement of \(w\) in the notation are not supplied. Consequently, \(T(d)w\) is retained as an unresolved transcription, without interpreting it as a product, subscript, or superscript.

The referenced Question 2 and the bibliography entries [4] and [5] are absent. Thus neither the geometric-realization requirement to be satisfied nor its assumptions, quantifiers, and categorical or cohomological comparison can be specified from this record. No separate conjecture or existence assertion is formulated here.

The supplied statement-verification note describes this sentence as the authors' proposed and partially established answer to a preceding geometric-realization question, rather than a separate open problem. The record labels its statement \texttt{unrecoverable} and supplies no clean statement. Its title uses ``Quasi-BPS,'' whereas the extracted sentence uses ``BPS category''; their intended relationship requires the missing context. This is a statement-only draft requiring human mathematical review.
\subsection{Short English statement}
The excerpt proposes a category as an answer to an omitted Question 2. Because that question, the category's definition, its parameters, and references [4] and [5] are missing, no precise standalone mathematical problem can be recovered. The supplied \texttt{partially\_solved} label is a curation label, not an independently verified mathematical conclusion.
\subsection{Sources}
\begin{itemize}
\item Local UnsolvedMath registry, \texttt{data/unsolvedmath/problems.json}, permanent ID \(30006270\), source code \texttt{catalogue entry}, title ``Reported Quasi-BPS Construction.'' The supplied record belongs to the Oberwolfach Reports Open Problems collection and records the proposed year as 2025. No individual proposer is supplied.
\item Supplied report citation: \emph{Enveloping Algebras and Geometric Representation Theory} (2025), Oberwolfach Reports, pp. 1271--1332, DOI 10.4171/owr/2025/25. \url{https://doi.org/10.4171/owr/2025/25}. The extracted text carries the report page marker 1311.
\item Supplied literature source, ``OWR 2025 report'': \url{https://ems.press/content/serial-article-files/51851}.
\item Supplied literature source, ``Quasi-BPS categories'': \url{https://arxiv.org/abs/2309.08425}.
\item Supplied literature source, ``Topological K-theory of quasi-BPS categories'': \url{https://arxiv.org/abs/2309.08432}.
\item Literature assessment dated exactly \texttt{2026-08-22}: supplied dated curation, not independently verified. It labels the problem partially solved and reports theorem-level categorical realizations of relevant K-theory and BPS cohomology by quasi-BPS categories in broad symmetric-quiver settings, while describing the requested general geometric and Hall-categorical interpretation as still having conjectural parts. The background identifies this assessment as literature triage derived from the supplied OWR curation and calls for independent verification of time-sensitive status and cited claims.
\end{itemize}
\end{document}
